\section{Verification and test suite}
\label{sec:tests}

The adaptations of Section~\ref{sec:adaptations} are guarded by an automated
\codeword{pytest} suite, \codeword{test_fss_consistency.py}. It serves two
purposes at once: it is a \emph{regression} suite that pins the behaviour of the
SAT code documented here, and a \emph{cross-method consistency} suite that checks
the SAT implementation against the two sibling backends in the repository --- the
FFT reference (\codeword{fss_FFT.py}) and the OpenCL implementation
(\codeword{fss_OCL.py}). Agreement between three independent code paths that
compute the same score is a strong correctness signal, and it is what lets the
performance rewrites of Sections~\ref{sec:filter}--\ref{sec:score} be trusted to
leave the result unchanged.

All tests run on a small set of reproducible Gaussian-bell fields built from a
fixed seed, so every value is deterministic. Two tolerances are used when
comparing methods, reflecting genuine algorithmic differences rather than bugs:

\begin{description}
  \item[$10^{-5}$] between SAT and OpenCL --- the \emph{same} algorithm, the only
        difference being \codeword{float32} arithmetic on the device versus
        \codeword{float64} on the host;
  \item[$10^{-3}$] between FFT and SAT --- the FFT convolution zero-pads at the
        domain boundary whereas the SAT filter edge-clamps
        (Section~\ref{sec:filter}), so the two agree in the interior but differ
        in a boundary frame.
\end{description}

The suite is run with

{\setstretch{1}\begin{lstlisting}[style=bash]
pytest test_fss_consistency.py -v
\end{lstlisting}}

\noindent and requires no arguments. The OpenCL tests detect whether a
\codeword{pyopencl} runtime is available and \emph{skip} themselves when it is
not, so the suite passes on a machine without a GPU/OpenCL stack.

\subsection{Coverage}
\label{sec:tests-coverage}

Table~\ref{tab:tests} groups the checks by theme. The right-hand column records
which of the three backends each theme exercises (S\,=\,SAT, F\,=\,FFT,
O\,=\,OpenCL).

\begin{table}[ht!]
\centering
\renewcommand{\arraystretch}{1.25}
\begin{tabularx}{\textwidth}{l X c}
\toprule
\textbf{Theme} & \textbf{What is verified} & \textbf{Backends} \\
\midrule
Cross-method
  & FFT, SAT and OpenCL scores agree within tolerance across thresholds
    $\times$ windows & S,F,O \\
Mathematical properties
  & perfect forecast $\to 1$, FSS $\in [0,1]$, NaN when no exceedances,
    monotonicity in window size, forecast/observation symmetry & S,F,O \\
Batch interfaces
  & \texttt{fss\_threshold}, \texttt{fss\_cumsum\_parallel} and
    \texttt{fss\_opencl\_arr} match the per-call single results & S,O \\
Regression guards
  & pinned reference scores for fixed (threshold, window) pairs, including
    every \texttt{threshold\_mode} & S \\
Threshold modes
  & \texttt{over}/\texttt{under}/\texttt{between}/\texttt{tolerance} preserve
    invariants and partition the domain correctly; the tolerance window is
    built from the forecast's own threshold on both bounds & S \\
\texttt{CWFSS}
  & R2 sequence, numerators / denominators / values / score and
    \texttt{bootstrap} are bit-identical to the reference, plus invariants
    and pinned values & S \\
Missing data
  & the masked path reduces \emph{exactly} to the clean path without NaNs, and
    preserves all invariants under injected NaNs (see
    Section~\ref{sec:tests-nan}) & S,F,O \\
\bottomrule
\end{tabularx}
\caption{Themes covered by \texttt{test\_fss\_consistency.py} and the backends
each one exercises (S\,=\,SAT, F\,=\,FFT, O\,=\,OpenCL).}
\label{tab:tests}
\end{table}

\subsection{The missing-data tests}
\label{sec:tests-nan}

Because the missing-data handling of Section~\ref{sec:nan} is the most
substantial addition, it carries the most checks, and they are mirrored across
all three backends (deterministic and ensemble for SAT and FFT; deterministic
for OpenCL, which has no ensemble entry point). The key guarantees are:

\begin{description}
  \item[Reduction to the clean path.] On NaN-free fields the masked score of
        Eq.~\eqref{eq:maskedscore} is compared against the ordinary fast path.
        For SAT and FFT this matches to machine precision ($\sim 10^{-9}$ or
        better), confirming that the valid count $C$ of
        Eq.~\eqref{eq:validcount} collapses to the constant window area $A$ when
        nothing is missing; for OpenCL it matches to the \codeword{float32}
        tolerance $10^{-5}$. A companion test verifies directly that
        $C \equiv A$ everywhere in the absence of NaNs.
  \item[Invariants under NaN.] With NaNs injected into both fields, the score
        stays in $[0,1]$, a perfect forecast (sharing the observation's NaNs)
        still scores $1$, the forecast/observation symmetry holds, and an
        all-missing field yields NaN rather than a spurious number.
  \item[Weighting semantics.] A hand-rolled weighted mean reproduces
        \codeword{_fss_score_masked}, pinning the per-window valid-point
        weighting of Eq.~\eqref{eq:maskedscore}.
  \item[Ensemble specifics.] A member that is entirely NaN does not change the
        score (the average is over valid members only), and the
        \codeword{nanpercentile} threshold path stays finite.
  \item[Cross-method agreement.] The FFT and OpenCL masked paths are checked
        against the SAT masked path. OpenCL also edge-clamps, so it agrees with
        SAT as closely as on clean data; FFT zero-pads, so it agrees within the
        interior at the usual $10^{-3}$ boundary tolerance.
  \item[Auto-dispatch.] Calling the public entry points
        (\codeword{fss_threshold}, \codeword{fourier_fss},
        \codeword{fourier_fss_eps}, \codeword{fss_opencl}) on a field that
        contains a NaN routes to the masked path automatically --- important for
        the FFT, where a raw NaN passed through \codeword{fftconvolve} would
        otherwise spread across the entire output.
\end{description}

\subsection{Result}
\label{sec:tests-result}

The suite currently comprises \textbf{308 tests}. With an OpenCL runtime present
all \textbf{308 pass}; on a machine without OpenCL the \textbf{60} device-backed
tests skip and the remaining \textbf{248 pass}. A representative run completes in
roughly six seconds.
